\section{Introduction}
\label{sec:introduction}

Peptide drug programs are gated by triage.
A typical campaign enters early discovery with thousands of in silico hits and a wet-lab budget that covers perhaps twenty.
The cost of a bad shortlist is not a wasted assay plate; it is six months of a medicinal chemistry program spent validating candidates that a better ranking would have deprioritized on day one.
Picking those twenty well requires weighing a portfolio of constraints simultaneously, since a peptide program that delivers on activity but fails on toxicity, stability, or manufacturability is not a lead \citep{yan2022review}.
Most computational pipelines still score one endpoint at a time \citep{veltri2018ampscanner,rathore2024toxinpred3,sharma2014hlp} and defer the multi-objective judgment to a spreadsheet.
The trade-offs are not imagined: highly active antimicrobial peptides tend to act through membrane disruption, a mechanism that is genuinely correlated with mammalian toxicity \citep{matsuzaki2009control}.
Single-endpoint winners are often multi-endpoint losers.

The textbook response is multi-objective optimization.
NSGA-II \citep{deb2002nsga2} and its descendants return a Pareto front: a diverse set of non-dominated compromises spread across the objective space \citep{emmerich2018tutorial}.
That output shape is a mismatch for the triage decision a program director actually makes.
Nobody needs a diverse frontier of 200 compromise peptides; they need the 20 candidates most likely to survive validation.
In practice, most teams fall back on hand-tuned weighted-sum scoring in a spreadsheet, because it at least returns a ranked list.
The weights are set by intuition, rarely revisited, and almost never audited against held-out data.

We asked whether an automated policy-search agent, given a frozen evaluation harness and a scored candidate pool, can learn a better ranking policy than the weighted sum a human team would write by hand.
The agent (Codex, GPT-5.4) reads the current ranking policy, edits the code, runs the harness, and keeps the change only if the score improves; otherwise it reverts.
Data, oracles, and metrics are read-only, so the agent can only improve by writing better ranking code, not by reshaping the test.
The loop design is borrowed from \citet{karpathy2026autoresearch} and applied in parallel to molecular model architectures in a companion study \citep{wijaya2026autonomousagentdiscoversmolecular}; here it is directed for the first time at peptide candidate triage.

The benchmark is public and modest in scope.
We assemble 3{,}554 antimicrobial peptides from DBAASP \citep{pirtskhalava2021dbaasp} with measured activity against \emph{E.~coli}, predicted toxicity from ToxinPred3 \citep{rathore2024toxinpred3}, predicted half-life from HLP \citep{sharma2014hlp}, and a rule-based manufacturability score.
The agent ran 100 experiments over roughly one laptop-day; 12 modifications passed the keep/discard gate.
Its final policy combines an oracle consensus voting stage with a learned reranker \citep{burges2010ranknet,ke2017lightgbm}, but the practical point for a biotech reader is that this structure was discovered automatically rather than prescribed by the authors.

Success is measured by top-20 enrichment (the fraction of the truly best candidates that land in the agent's top-20 shortlist), averaged across three independent oracle definitions, with NDCG as a secondary metric and bootstrap confidence intervals over 10 random data splits.
The agent-derived policy captures 65\% of the best candidates in its top-20, compared to 44\% for NSGA-II and 61\% for an equal-weight scalarization baseline.
The margin over NSGA-II is large and its bootstrap confidence intervals are cleanly separated; the margin over weighted-sum is smaller, and we are honest about it: single-split confidence intervals overlap, while the cross-split mean separates consistently ($0.650 \pm 0.034$ vs.\ $0.585 \pm 0.050$ across 10 splits).
All claims are on a public AMP benchmark; we are not making a clinical claim.

This paper contributes the following, framed in terms of what a peptide-program reader gets:
\begin{enumerate}[nosep]
\item \textbf{A drop-in triage layer for peptide programs.}
The released ranking policy, evaluation harness, and loop infrastructure are designed to replace a weighted-scoring spreadsheet; swapping in an internal peptide library and in-house endpoint models requires changing only the data pipeline.

\item \textbf{A public multi-objective AMP benchmark.}
A set of 3{,}554 peptides scored across activity, toxicity, stability, and manufacturability, evaluated under three mathematically distinct oracle definitions, is released so other groups can benchmark their own triage methods on a common target.

\item \textbf{Evidence that automated policy search beats standard baselines on triage.}
On this benchmark, the agent-derived policy outperforms NSGA-II and equal-weight scalarization on top-20 enrichment, with cross-split consistency across 10 random splits.

\item \textbf{The agent earned its keep beyond prompt-wrapped random search.}
The final policy contains structural elements (oracle consensus voting plus a learned reranker) that neither the authors pre-specified nor a 1{,}000-sample random weight search could reach (0.650 vs.\ 0.611 top-20 enrichment), indicating that the loop is discovering policy structure, not just retuning weights.
\end{enumerate}
